%%% Šablona pro jednoduchý soubor formátu PDF/A, jako treba samostatný abstrakt práce.

\documentclass[12pt]{report}

\usepackage[a4paper, hmargin=1in, vmargin=1in]{geometry}
\usepackage[a-2u]{pdfx}
\usepackage[T1]{fontenc}
\usepackage[utf8]{inputenc}
\usepackage[czech,english]{babel}

\usepackage{lmodern}
\usepackage{textcomp}
\usepackage{fontspec}
\usepackage{polyglossia}

\setmainlanguage{english}
\setotherlanguage{czech}
\begin{document}

%% Nezapomeňte upravit abstrakt.xmpdata.

Tato práce zkoumá využití novelty v evolučních algoritmech při řešení úloh zpětnovazebního učení. Porovnává selekci založenou na novelty se selekcí založenou na fitness, jejich variantami využívajícími archivy a dynamickými kombinacemi novelty a fitness. Porovnání probíhá na prostředích CartPole a Lunar Lander z knihovny Gymnasium. Pro každé prostředí definujeme interpretovatelnou behaviorální charakteristiku a analyzujeme vliv jednotlivých metod práce s objektivní funkcí na dosaženou fitness a diverzitu populace.
Výsledky ukazují, že metody založené na fitness představují nejsilnější výchozí přístup z hlediska dosaženého výkonu. Využití novelty zvyšuje diverzitu populace, avšak její přínos závisí na dynamice konkrétní úlohy.
Práce dále diskutuje citlivost výsledků algoritmů založených na novelty na volbu behaviorálního prostoru a vzájemné interakce mezi prostředím a evolučním algoritmem.

\end{document}
